% 统计显著性检验表
\documentclass{article}
\usepackage[utf8]{inputenc}
\usepackage{booktabs}
\usepackage{amsmath}
\usepackage{caption}
\captionsetup[table]{labelformat=empty}
\renewcommand{\arraystretch}{1.2}

\begin{document}

\begin{table}[htbp]
\centering
\caption{Statistical Significance of Performance Gains (15-fold LOSO, paired tests)}
\label{tab:significance}
\begin{tabular}{lcccccc}
\toprule
\textbf{Comparison} &
\textbf{MSE (A)} &
\textbf{MSE (B)} &
\textbf{$\Delta$ MSE} &
\textbf{Paired $t$} &
\textbf{$p$ (t-test)} &
\textbf{$p$ (Wilcoxon)} \\
\midrule
vs.\ Baseline A     & 0.1818 & 0.1417 & $-0.0401$ & 1.67 & 0.117  & \textbf{0.018}$^{*}$ \\
vs.\ Baseline B     & 0.1903 & 0.1417 & $-0.0487$ & 2.81 & \textbf{0.014}$^{*}$ & \textbf{0.008}$^{**}$ \\
vs.\ w/o Dual Gating & 0.1458 & 0.1417 & $-0.0041$ & 0.52 & 0.610  & 0.762 \\
vs.\ w/o TCM Prior  & 0.1516 & 0.1417 & $-0.0099$ & 1.09 & 0.294  & 0.359 \\
\bottomrule
\end{tabular}
\vspace{0.5em}

\footnotesize{
$^{*}p < 0.05$, \quad $^{**}p < 0.01$. \quad
Wilcoxon signed-rank test is preferred due to non-normal MSE distribution (outlier in S7).
The proposed model significantly outperforms both dynamic-only baselines.
Ablation variants (w/o Dual Gating, w/o TCM Prior) do not reach significance,
consistent with their smaller MSE gaps.
}
\end{table}

\end{document}
